\section{Machine-control co-design and an implementable design loop}
\label{sec:codesign}

Machine design and current allocation share the same quadratic forms.  A
change in \(L_d-L_q\) changes the torque surface, its loss-optimal eigenvector,
the voltage ellipse, and the field-weakening trajectory simultaneously.  It is
therefore artificial to select a machine first and ask only afterward how it
should be controlled.  The coupled object is
\begin{equation}
 \min_{\vect g}\ J(\vect g,\vect i_1^\star,\ldots,\vect i_N^\star)
 \quad\text{subject to}\quad
 \vect i_k^\star\in\argmax_{\vect i}\{M_k:\ P_k\leq P_{\max},
 U_k\leq U_{\max}\},
 \label{eq:machine-control-codesign}
\end{equation}
with \(\boldsymbol\theta=\mathcal M(\vect g)\).  Part I supplies the inner
solution; Part II turns its optimal value into outer constraints.

\subsection{PSM points and EESM trajectories in one picture}

A PSM has one fixed excitation coordinate \(\psi_{\mathrm{pm}}\).  An EESM
generates a family of effective PSM slices
\begin{equation}
 \psi_f(i_{\mathrm{exc}})=L_{\mathrm{exc}}i_{\mathrm{exc}},
 \qquad
 M=k_Ti_q(\Delta L i_d+\psi_f).
 \label{eq:effective-pm-trajectory}
\end{equation}
As \(i_{\mathrm{exc}}\) varies, the EESM traces an effective-PM-flux path
through PSM-like operating geometries.  This extra degree of freedom can move
the voltage ellipse and trade stator current against field current across
speed.  Its price is the additional term
\(R_{\mathrm{exc}}i_{\mathrm{exc}}^2\), an excitation converter, rotor thermal
loading, and field dynamics.  Thus an EESM does not dominate every PSM point
for free; it buys a controllable trajectory through design/performance space.

\subsection{Outer optimizer, inner geometric oracle}

For each proposed \(\boldsymbol\theta\), an implementable oracle is:
\begin{enumerate}
 \item reject immediately if the closed loss-only upper bound is below the
       requested torque;
 \item for each operating point, search the support slope \(\mu\) and solve a
       symmetric \(3\times3\) generalized eigenproblem;
 \item scale the selected direction by the active loss or voltage radius;
 \item return currents, active constraints, \(P_{\min}\), \(M_{\max}\), and
       signed margins;
 \item for nonlinear maps, update the local affine model and verify the final
       point in the original flux, voltage, and thermal equations.
\end{enumerate}
This regular small computation is suitable for screening analytical designs,
initializing numerical optimization, training a surrogate on physically
meaningful outputs, or filtering candidates before \gls{fea}.  It also matches
the embedded controller: the offline design tool and online current-reference
generator can use the same tested value-function kernel with different
parameter sources.

For an \gls{mcu}, the most useful result is the bounded scalar support search,
closed inactive-voltage branch, and explicit feasibility margin.  For an
\gls{fpga}, a fixed-sweep Jacobi eigensolver or tabulated support frontier gives
regular data flow.  The inverse design extension does not increase online
state dimension; it clarifies which parameter combinations and uncertainty
ranges the kernel must support.

\subsection{What remains exact, local, or open}

The EESM capability circle, its sensitivities, and the nonsalient PSM flux
bound are exact only for constant linear parameters with voltage inactive.
The characteristic-current upper bound additionally uses the ideal
flux-cancellation interpretation.  Voltage-aware design boundaries are exact
for the stated linear steady-state model but generally implicit.  Local affine
flux maps retain the value-function architecture under saturation, while the
global circle and constant spectral descriptors become parameter paths.

Open mathematical questions include conditions for convexity or star-convexity
of useful low-dimensional design slices, certified error bounds when a
surrogate approximates \(\mathcal M\), nonsmooth sensitivities at active-set
changes and repeated eigenvalues, and contraction certificates for the
nonlinear outer flux update.  A further practical question is how few support
values of \(\widetilde H-\mu\widetilde Q\) suffice to fingerprint voltage-aware
capability without missing a narrow torque-speed constraint.

